\section{Related Work}
\label{sec:related}

\paragraph{Process models.}
Carrier and Spafford~\cite{carrier2003getting} describe \emph{when} digital
evidence is collected; ours describes \emph{how} it is reached. The
hypothesis-based framework~\cite{carrier2006hypothesis} presupposes an
investigator knows where to look; our taxonomy names the five address spaces
where evidence can live. Beebe and Clark~\cite{beebe2005hierarchical},
DFRWS~\cite{palmer2001roadmap}, and NIST SP 800-86~\cite{kent2006sp80086}
provide procedural guidance that our taxonomy informs but does not replace.

\paragraph{Data-abstraction layer models.}
Carrier's data-abstraction model~\cite{carrier2005file} asks \emph{which
abstraction level} an artifact inhabits (physical $\to$ volume $\to$
filesystem $\to$ file). Our taxonomy asks \emph{which addressing mechanism}
reaches bytes at any level. The two axes jointly span a 2-D classification
space; TSK instantiates the (\src{P}, filesystem) cell. TSK has no
analog for \src{M}, \src{Q}, or \src{C}; our taxonomy generalizes
Carrier's layered insight to four addressing schemes that postdate TSK's
design. Garfinkel's DFXML work~\cite{garfinkel2009automating} provides
early evidence that even within \src{P} alone, absent a uniform addressing
model, ad-hoc adapters proliferate.

\paragraph{Representation ontologies.}
CASE/UCO~\cite{caseuco2018,casey2017advancing} represents evidence
\emph{after} parsing. Our taxonomy governs the architecture by which
evidence is reached and parsed. A CASE/UCO record can carry a
\texttt{source-type} attribute from our taxonomy, providing provenance that
the representation alone does not encode.

\paragraph{Source-specific tooling.}
TSK/Autopsy~\cite{carrier2005file} (\src{P}), Volatility~\cite{ligh2014artmemory}
(\src{M}), and Plaso~\cite{plaso} (\src{L}) each excel within one primitive.
Our taxonomy provides the formal basis for composing them: all produce
$\Bstar$, so a parser-of-record can accept input from any without
modification.

\paragraph{Live response.}
Cohen, Bilby, and Garfinkel~\cite{cohen2011distributed} argued that
live-query collection is the realistic alternative to full disk imaging at
enterprise scale but did not formalize what live-query evidence \emph{is}.
Our \src{Q} primitive and Proposition~\ref{prop:q-durability} supply that
formalization: evidence produced by query execution, once transmitted,
is durable against on-host adversarial action.

\paragraph{Content-addressed foundations.}
in-toto~\cite{torresarias2019intoto}, Sigstore~\cite{newman2022sigstore},
OCI~\cite{ociimagespec}, and IPFS~\cite{benet2014ipfs} established
content-addressing as critical for supply-chain integrity. We incorporate
\src{C} into the forensic taxonomy as a first-class peer, with explicit
forensic semantics (Definition~\ref{def:hash-pivot}) and an intrinsic
durability characterization (Proposition~\ref{prop:c-durability}).

\paragraph{Closest prior work.}
Garfinkel's 2010 agenda~\cite{garfinkel2010digital} identified cross-source
correlation as an open problem and called for ``a model that allows multiple
sources of evidence to be combined in principled ways.'' Our taxonomy
supplies the missing piece: the navigation-primitive tag is the formally
identified provenance attribute a correlation engine needs to reason about
trust and apply source-appropriate durability profiles.
